
quasi-compacts et quasi-s\'epar\'es, \`a morphismes de transition affines, l'application canonique}
@@M0@@
\emph{est bijective.}

\emph{Par cons\'equent} (EGA IV 8.14.2), \emph{si le foncteur de Picard est repr\'esent\'e par un @@M1@@-sch\'ema @@M2@@, alors @@M3@@ est localement de pr\'esentation finie.}

En effet, un \'el\'ement @@M4@@ de @@M5@@ est repr\'esent\'e par un faisceau inversible @@M6@@ sur un @@M7@@ o\`u @@M8@@ est un morphisme f.p.p.f. Posons @@M9@@. L'image @@M10@@ de @@M11@@ dans @@M12@@ est repr\'esent\'ee par le faisceau inversible @@M13@@ sur @@M14@@, image inverse de @@M15@@, o\`u @@M16@@.

Si @@M17@@ est l'\'el\'ement neutre de @@M18@@, il existe un morphisme f.p.p.f. @@M19@@ tel que l'image inverse de @@M20@@ sur @@M21@@ soit @@M22@@. En vertu de (EGA IV 8.8.2, 8.10.5 (vi), 11.2.6 (ii), 8.5.2 (i), et 8.5.2.4), il existe un @@M23@@ et un morphisme f.p.p.f. @@M24@@ o\`u @@M25@@ tel que l'image inverse de @@M26@@ sur @@M27@@ soit @@M28@@. L'image de @@M29@@ est donc d\'ej\`a l'\'el\'ement neutre dans @@M30@@. Par suite 3.1.1 est injective.

La surjectivit\'e de 3.1.1 se d\'emontre d'une fa\c{c}on analogue.

\textbf{Proposition 3.2.} \emph{Soient @@M31@@ un sch\'ema noeth\'erien et @@M32@@ un @@M33@@-sch\'ema propre. On suppose que le foncteur de Picard est repr\'esent\'e par un @@M34@@-sch\'ema @@M35@@. Alors}~:
\begin{enumerate}[label=(\roman*)]
\item \emph{@@M36@@ est localement de type fini sur @@M37@@.}
